%!TEX TS-program = xelatex
%!TEX encoding = UTF-8 Unicode
\documentclass[10pt]{article}
\title{Java vs TypeScript --- A cheatsheet}
\author{Mathieu N\'eron}
\usepackage{cheatsheet-style}

\begin{document}
\cheatsheetTitle{Java vs TypeScript}{Mathieu N\'eron}{2026}

\noindent\small This cheatsheet anchors TypeScript on constructs you already know in
Java. Left column: Java. Right column: TypeScript (5.x). The biggest mental shifts:
\textbf{structural typing} (shape, not name), \textbf{single-threaded event loop}
with \code{async}/\code{await}, \textbf{union/literal types}, and a vast \textbf{type
system that erases at runtime}.

%===============================================================================
\cheatsection{Hello world \& program entry}
%===============================================================================

\begin{construct}{Program entry point}
\noindent
\begin{minipage}[t]{0.485\linewidth}
\begin{lstlisting}[style=java, title=Java]
public class Hello {
  public static void main(String[] args) {
    System.out.println("Hello, " + args[0]);
  }
}
// javac Hello.java && java Hello world
\end{lstlisting}
\end{minipage}\hfill
\begin{minipage}[t]{0.485\linewidth}
\begin{lstlisting}[style=ts, title=TypeScript]
// hello.ts (Node 20+ / Bun / Deno)
const name: string = process.argv[2];
console.log(`Hello, ${name}`);

// $ tsx hello.ts world      (or: tsc && node hello.js)
// In a module, top-level code just runs at load.
\end{lstlisting}
\end{minipage}
\end{construct}

\begin{construct}{File / module / package layout}
\noindent
\begin{minipage}[t]{0.485\linewidth}
\begin{lstlisting}[style=java, title=Java]
// One public class per file.
package com.acme;
// pom.xml / build.gradle declares deps.
\end{lstlisting}
\end{minipage}\hfill
\begin{minipage}[t]{0.485\linewidth}
\begin{lstlisting}[style=ts, title=TypeScript]
// Each .ts file is a module if it has import/export.
// package.json + node_modules. Entry: "main" or "exports".
// Scripts via "scripts": { ... }; pnpm/npm/bun.
// tsconfig.json controls compiler.
\end{lstlisting}
\end{minipage}
\end{construct}

%===============================================================================
\cheatsection{Variables, constants, type inference}
%===============================================================================

\begin{construct}{Local variables}
\noindent
\begin{minipage}[t]{0.485\linewidth}
\begin{lstlisting}[style=java, title=Java]
int x = 5;
final String s = "hi";
var list = new ArrayList<Integer>();   // Java 10+
\end{lstlisting}
\end{minipage}\hfill
\begin{minipage}[t]{0.485\linewidth}
\begin{lstlisting}[style=ts, title=TypeScript]
let x: number = 5;            // mutable
const s: string = "hi";       // immutable BINDING
const list: number[] = [];    // type can be inferred
const arr = [1, 2, 3];        // inferred number[]
// `var` exists but avoid it — function-scoped, hoisted.
\end{lstlisting}
\end{minipage}
\end{construct}

\begin{construct}{Const != deeply immutable}
\noindent
\begin{minipage}[t]{0.485\linewidth}
\begin{lstlisting}[style=java, title=Java]
final List<Integer> xs = new ArrayList<>(List.of(1,2));
xs.add(3);                  // OK — final binding only
List<Integer> immut = List.of(1, 2);
immut.add(3);               // throws UnsupportedOp
\end{lstlisting}
\end{minipage}\hfill
\begin{minipage}[t]{0.485\linewidth}
\begin{lstlisting}[style=ts, title=TypeScript]
const xs = [1, 2];
xs.push(3);                  // OK — const binding only
const immut: readonly number[] = [1, 2];
immut.push(3);               // type error
const tuple = [1, 2] as const;   // readonly tuple [1,2]
\end{lstlisting}
\end{minipage}
\end{construct}

%===============================================================================
\cheatsection{Primitive types \& conversions}
%===============================================================================

\begin{construct}{Numeric \& boolean types}
\noindent
\begin{minipage}[t]{0.485\linewidth}
\begin{lstlisting}[style=java, title=Java]
byte/short/int/long; float/double;
int n = Integer.parseInt("42");
double x = Double.parseDouble("1.5");
String t = String.valueOf(42);
\end{lstlisting}
\end{minipage}\hfill
\begin{minipage}[t]{0.485\linewidth}
\begin{lstlisting}[style=ts, title=TypeScript]
// number = 64-bit IEEE float (no separate int).
// bigint = arbitrary precision (literal: 42n).
const n: number = parseInt("42", 10);
const x: number = parseFloat("1.5");
const t: string = String(42);
const big: bigint = 42n;
\end{lstlisting}
\end{minipage}
\end{construct}

\begin{construct}{Strings (template literals)}
\noindent
\begin{minipage}[t]{0.485\linewidth}
\begin{lstlisting}[style=java, title=Java]
String name = "world";
String s1 = "Hello, " + name;
String s2 = String.format("%s = %.2f", "pi", 3.14);
String body = """
    Hello,
    World
    """;
\end{lstlisting}
\end{minipage}\hfill
\begin{minipage}[t]{0.485\linewidth}
\begin{lstlisting}[style=ts, title=TypeScript]
const name = "world";
const s1 = `Hello, ${name}`;
const s2 = `${"pi"} = ${(3.14).toFixed(2)}`;
const body = `
    Hello,
    World
    `;
// Tagged templates: fn`a ${x} b` -> fn(strings, x)
\end{lstlisting}
\end{minipage}
\end{construct}

%===============================================================================
\cheatsection{Control flow}
%===============================================================================

\begin{construct}{If / else / ternary / nullish}
\noindent
\begin{minipage}[t]{0.485\linewidth}
\begin{lstlisting}[style=java, title=Java]
if (x > 0) { ... } else if (x == 0) { ... } else { ... }
int sign = x > 0 ? 1 : (x < 0 ? -1 : 0);

String name = (s != null) ? s : "default";
\end{lstlisting}
\end{minipage}\hfill
\begin{minipage}[t]{0.485\linewidth}
\begin{lstlisting}[style=ts, title=TypeScript]
if (x > 0) { ... } else if (x === 0) { ... } else { ... }
const sign = x > 0 ? 1 : (x < 0 ? -1 : 0);

const name = s ?? "default";        // nullish-coalescing
const name2 = s || "default";       // FALSY (incl. "", 0)
const len = obj?.user?.name?.length;// optional chaining
\end{lstlisting}
\end{minipage}
\end{construct}

\begin{construct}{Switch / discriminated unions}
\noindent
\begin{minipage}[t]{0.485\linewidth}
\begin{lstlisting}[style=java, title=Java]
String label = switch (shape) {
  case Circle c       -> "circle r=" + c.r();
  case Rectangle(var w, var h) -> "rect " + w + "x" + h;
  default             -> "unknown";
};
\end{lstlisting}
\end{minipage}\hfill
\begin{minipage}[t]{0.485\linewidth}
\begin{lstlisting}[style=ts, title=TypeScript]
type Shape =
  | { kind: "circle"; r: number }
  | { kind: "rect"; w: number; h: number };

function label(s: Shape): string {
  switch (s.kind) {
    case "circle": return `circle r=${s.r}`;
    case "rect":   return `rect ${s.w}x${s.h}`;
  }
  // Exhaustiveness: TS narrows `s` to never above.
}
\end{lstlisting}
\end{minipage}
\end{construct}

%===============================================================================
\cheatsection{Loops}
%===============================================================================

\begin{construct}{Counting / iteration}
\noindent
\begin{minipage}[t]{0.485\linewidth}
\begin{lstlisting}[style=java, title=Java]
for (int i = 0; i < 10; i++) { ... }
for (var x : xs) { ... }
for (var e : map.entrySet()) {
    System.out.println(e.getKey() + "=" + e.getValue());
}
\end{lstlisting}
\end{minipage}\hfill
\begin{minipage}[t]{0.485\linewidth}
\begin{lstlisting}[style=ts, title=TypeScript]
for (let i = 0; i < 10; i++) { ... }

for (const x of xs)        { ... }      // values
for (const [i, x] of xs.entries()) { ... }
for (const [k, v] of m) { ... }         // Map
for (const k in obj) { ... }            // KEYS (rarely)

xs.forEach((x, i) => { ... });          // higher-order
\end{lstlisting}
\end{minipage}
\end{construct}

\begin{construct}{Async iteration}
\noindent
\begin{minipage}[t]{0.485\linewidth}
\begin{lstlisting}[style=java, title=Java]
// Reactive stack or a custom Iterator. No syntax sugar.
\end{lstlisting}
\end{minipage}\hfill
\begin{minipage}[t]{0.485\linewidth}
\begin{lstlisting}[style=ts, title=TypeScript]
for await (const chunk of stream) {     // AsyncIterable
    process(chunk);
}
\end{lstlisting}
\end{minipage}
\end{construct}

%===============================================================================
\cheatsection{Collections \& data structures}
%===============================================================================

JavaScript's built-ins are: \code{Array}, \code{Map}, \code{Set},
\code{WeakMap}, \code{WeakSet}, plus typed-array buffers. Both \code{Map} and
\code{Set} preserve insertion order --- so they double as
\code{LinkedHashMap}/\code{LinkedHashSet} out of the box. For sorted maps,
priority queues, and deques you'll reach for npm libraries (\code{mnemonist},
\code{js-sdsl}, or \code{sorted-btree}).

\cheatsub{Quick reference}

\noindent
\begin{tabularx}{\linewidth}{@{}>{\ttfamily\small}l >{\small}l X@{}}
\toprule
\rowcolor{accentDark}
\refhead{\normalfont Java} & \refhead{Behavior} & \refhead{TypeScript} \\
\midrule
ArrayList<T>        & dynamic array & \code{T[]} (\code{Array<T>}) \\
LinkedList<T>       & doubly-linked & \code{mnemonist/linked-list}; or just use Array \\
int[] / Object[]    & fixed-size array & \code{Int32Array(N)} typed array; \code{Array(N)} (sparse) \\
HashSet<T>          & unordered set & \code{Set<T>} (note: insertion-ORDERED) \\
LinkedHashSet<T>    & insertion-ordered set & \code{Set<T>} (it already is) \\
TreeSet<T>          & sorted set & \code{js-sdsl} \code{OrderedSet}; \code{sorted-btree} \\
EnumSet             & bit-set over enum & \code{Set<MyEnum>} \\
HashMap<K,V>        & unordered map & \code{Map<K,V>} (insertion-ORDERED); object-as-dict for string keys \\
LinkedHashMap<K,V>  & insertion-ordered & \code{Map<K,V>} (it already is) \\
TreeMap<K,V>        & sorted map & \code{js-sdsl} \code{OrderedMap}; \code{sorted-btree} \\
IdentityHashMap     & identity-keyed & \code{Map<K,V>}: object keys ARE compared by reference \\
EnumMap             & enum-keyed map & \code{Map} or \code{Record<MyEnum, V>} \\
WeakHashMap         & weak keys & \code{WeakMap<K,V>} (keys must be objects) \\
ArrayDeque          & double-ended queue & \code{mnemonist/deque}; Array (shift is O(n)) \\
ArrayDeque (queue)  & FIFO & \code{mnemonist/deque} or \code{Array.shift} (slow) \\
ArrayDeque (stack)  & LIFO & \code{Array.push/pop} \\
PriorityQueue<T>    & min-heap & \code{mnemonist/heap}, \code{js-sdsl} \code{PriorityQueue} \\
ConcurrentHashMap   & thread-safe map & N/A --- single-threaded event loop; cross-worker via \code{postMessage} \\
BlockingQueue       & blocking FIFO & async generator + \code{await}; or \code{node:stream} \\
Counter             & multiset / tally & \code{Map<K,number>}; \code{Object.groupBy} (ES2024) \\
List.copyOf, Map.of & immutable view & \code{readonly T[]}, \code{ReadonlyMap<K,V>} (compile-time only); \code{Object.freeze} for runtime \\
\bottomrule
\end{tabularx}

\bigskip

\cheatsub{Arrays, sets, maps (the natives)}

\begin{construct}{Array (ArrayList analog)}
\noindent
\begin{minipage}[t]{0.485\linewidth}
\begin{lstlisting}[style=java, title=Java]
List<Integer> xs = new ArrayList<>(List.of(1,2,3));
xs.add(4); xs.get(0); xs.size(); xs.remove(0);
xs.subList(1, 3); Collections.reverse(xs);
\end{lstlisting}
\end{minipage}\hfill
\begin{minipage}[t]{0.485\linewidth}
\begin{lstlisting}[style=ts, title=TypeScript]
const xs: number[] = [1, 2, 3];
xs.push(4); xs[0]; xs.length; xs.splice(0, 1);
xs.slice(1, 3);              // copy [1,3)
xs.toReversed();             // ES2023 (or [...xs].reverse())

// Higher-order:
xs.filter(x => x > 0).map(x => x*2);
xs.reduce((a, x) => a + x, 0);
xs.find(x => x === 2); xs.includes(2); xs.at(-1);

// Typed arrays (ArrayBuffer-backed, no boxing):
const buf = new Int32Array(1024);
\end{lstlisting}
\end{minipage}
\end{construct}

\begin{construct}{Map (LinkedHashMap by default)}
\noindent
\begin{minipage}[t]{0.485\linewidth}
\begin{lstlisting}[style=java, title=Java]
Map<String,Integer> m = new HashMap<>();
m.put("a", 1); m.getOrDefault("b", 0);
m.computeIfAbsent("c", k -> heavy());
m.merge("a", 1, Integer::sum);
\end{lstlisting}
\end{minipage}\hfill
\begin{minipage}[t]{0.485\linewidth}
\begin{lstlisting}[style=ts, title=TypeScript]
const m = new Map<string, number>();
m.set("a", 1); m.get("a") ?? 0;
m.has("a"); m.delete("a"); m.size;

// computeIfAbsent idiom:
let v = m.get("c") ?? (() => { const x = heavy(); m.set("c", x); return x; })();

// merge / tally:
m.set("a", (m.get("a") ?? 0) + 1);

// Iteration order = insertion order (always).
for (const [k, v] of m) { ... }

// Map vs object: Map keys can be ANY type. Object keys are
// always strings (or symbols). Map has .size; objects don't.
const r: Record<string, number> = { a: 1 };
\end{lstlisting}
\end{minipage}
\end{construct}

\begin{construct}{Set (LinkedHashSet by default)}
\noindent
\begin{minipage}[t]{0.485\linewidth}
\begin{lstlisting}[style=java, title=Java]
Set<String> hs = new HashSet<>();
Set<String> ls = new LinkedHashSet<>();
Set<String> ts = new TreeSet<>();
\end{lstlisting}
\end{minipage}\hfill
\begin{minipage}[t]{0.485\linewidth}
\begin{lstlisting}[style=ts, title=TypeScript]
const s = new Set<string>(["a", "b"]);
s.add("c"); s.has("a"); s.delete("b"); s.size;

// Set ops (ES2025):
s.union(t); s.intersection(t); s.difference(t);
s.symmetricDifference(t); s.isSubsetOf(t);

// Pre-ES2025 polyfill:
const u = new Set([...s, ...t]);                  // union
const i = new Set([...s].filter(x => t.has(x)));  // inter
\end{lstlisting}
\end{minipage}
\end{construct}

\cheatsub{Weak collections}

\begin{construct}{WeakMap / WeakSet}
\noindent
\begin{minipage}[t]{0.485\linewidth}
\begin{lstlisting}[style=java, title=Java]
Map<Key,V> wk = new WeakHashMap<>();
\end{lstlisting}
\end{minipage}\hfill
\begin{minipage}[t]{0.485\linewidth}
\begin{lstlisting}[style=ts, title=TypeScript]
const wm = new WeakMap<object, V>();      // keys MUST be objects
wm.set(keyObj, value);
wm.get(keyObj); wm.has(keyObj); wm.delete(keyObj);
// No iteration, no .size --- entries vanish silently when
// the key is GC'd. Useful for attaching metadata to objects
// without preventing collection.

const ws = new WeakSet<object>();
\end{lstlisting}
\end{minipage}
\end{construct}

\cheatsub{Sorted / priority structures (use libraries)}

\begin{construct}{TreeMap / TreeSet (sorted)}
\noindent
\begin{minipage}[t]{0.485\linewidth}
\begin{lstlisting}[style=java, title=Java]
TreeMap<String,Integer> tm = new TreeMap<>();
tm.firstKey(); tm.headMap("m"); tm.floorKey("g");

TreeSet<Integer> ts = new TreeSet<>();
ts.first(); ts.subSet(1, 10);
\end{lstlisting}
\end{minipage}\hfill
\begin{minipage}[t]{0.485\linewidth}
\begin{lstlisting}[style=ts, title=TypeScript]
// No stdlib. Pick a library:
import { OrderedMap, OrderedSet } from "js-sdsl";
const tm = new OrderedMap<string, number>();
tm.setElement("b", 2); tm.setElement("a", 1);
tm.front()!.first;       // smallest key

// Or sorted-btree:
import BTree from "sorted-btree";
const t = new BTree<string, number>();
\end{lstlisting}
\end{minipage}
\end{construct}

\begin{construct}{PriorityQueue (heap)}
\noindent
\begin{minipage}[t]{0.485\linewidth}
\begin{lstlisting}[style=java, title=Java]
PriorityQueue<Job> pq = new PriorityQueue<>(
    Comparator.comparingInt(Job::priority));
pq.offer(j); Job next = pq.poll();
\end{lstlisting}
\end{minipage}\hfill
\begin{minipage}[t]{0.485\linewidth}
\begin{lstlisting}[style=ts, title=TypeScript]
import Heap from "mnemonist/heap";
const pq = new Heap<Job>((a, b) => a.priority - b.priority);
pq.push(j); const next = pq.pop();

// Or js-sdsl PriorityQueue, or hand-roll an array-based
// binary heap (~30 lines) if you want zero deps.
\end{lstlisting}
\end{minipage}
\end{construct}

\cheatsub{Deque, queue, stack}

\begin{construct}{Stack and queue patterns}
\noindent
\begin{minipage}[t]{0.485\linewidth}
\begin{lstlisting}[style=java, title=Java]
Deque<Integer> st = new ArrayDeque<>();    // stack
st.push(1); st.peek(); st.pop();

Deque<Integer> q = new ArrayDeque<>();     // FIFO
q.offer(1); q.poll();
\end{lstlisting}
\end{minipage}\hfill
\begin{minipage}[t]{0.485\linewidth}
\begin{lstlisting}[style=ts, title=TypeScript]
// Stack: just an array.
const st: number[] = [];
st.push(1); const top = st.at(-1); const v = st.pop();

// Queue: Array.shift() works but is O(n) per call.
// For hot paths use mnemonist/deque (O(1) both ends):
import Deque from "mnemonist/deque";
const q = new Deque<number>();
q.push(1); q.shift();         // O(1)
\end{lstlisting}
\end{minipage}
\end{construct}

\cheatsub{Specialized}

\begin{construct}{Counter / groupBy / immutable views}
\noindent
\begin{minipage}[t]{0.485\linewidth}
\begin{lstlisting}[style=java, title=Java]
Map<String,Long> tally = words.stream()
    .collect(Collectors.groupingBy(w -> w, Collectors.counting()));
List<Integer> immut = List.copyOf(xs);
\end{lstlisting}
\end{minipage}\hfill
\begin{minipage}[t]{0.485\linewidth}
\begin{lstlisting}[style=ts, title=TypeScript]
// Tally:
const counts = new Map<string, number>();
for (const w of words) counts.set(w, (counts.get(w) ?? 0) + 1);

// ES2024 group-by (returns plain object / Map):
const byLen = Object.groupBy(words, w => w.length);
const byLenMap = Map.groupBy(words, w => w.length);

// Compile-time read-only (erased at runtime):
const xs: readonly number[] = [1,2,3];
const m: ReadonlyMap<string, number> = new Map();

// Runtime freeze (shallow):
const frozen = Object.freeze({ a: 1 });
\end{lstlisting}
\end{minipage}
\end{construct}

\begin{construct}{Tuples \& destructuring}
\noindent
\begin{minipage}[t]{0.485\linewidth}
\begin{lstlisting}[style=java, title=Java]
record Pair<A,B>(A first, B second) {}
Pair<Integer,String> p = new Pair<>(1, "x");
int a = p.first(); String b = p.second();
\end{lstlisting}
\end{minipage}\hfill
\begin{minipage}[t]{0.485\linewidth}
\begin{lstlisting}[style=ts, title=TypeScript]
const pair: [number, string] = [1, "x"];
const [a, b] = pair;                     // tuple destructure
const { x, y } = point;                  // object destructure
const { x: px, y: py } = point;          // rename
const [head, ...rest] = xs;              // rest pattern
const ro = [1, 2] as const;              // readonly tuple
\end{lstlisting}
\end{minipage}
\end{construct}

\begin{construct}{Mutability defaults}
\noindent
\begin{minipage}[t]{\linewidth}
\small
\textbf{Java}: \code{ArrayList}, \code{HashMap}, \code{HashSet} mutable; \code{List.of/Set.of/Map.of} are runtime-immutable.\par
\textbf{TypeScript}: All native collections are mutable at runtime. \code{readonly} / \code{Readonly<T>} / \code{ReadonlyMap} / \code{ReadonlySet} are \emph{compile-time only} --- they erase at runtime. For real runtime immutability use \code{Object.freeze} (shallow), \code{structuredClone} for deep copies, or libraries like \code{immer} / \code{immutable.js}.
\end{minipage}
\end{construct}

%===============================================================================
\cheatsection{Functions}
%===============================================================================

\begin{construct}{Definition, default args, varargs}
\noindent
\begin{minipage}[t]{0.485\linewidth}
\begin{lstlisting}[style=java, title=Java]
int add(int a, int b) { return a + b; }
int add(int a) { return add(a, 1); }    // overload
int sum(int... xs) { ... }
\end{lstlisting}
\end{minipage}\hfill
\begin{minipage}[t]{0.485\linewidth}
\begin{lstlisting}[style=ts, title=TypeScript]
function add(a: number, b: number = 1): number {
    return a + b;
}
function sum(...xs: number[]): number {
    return xs.reduce((t, x) => t + x, 0);
}
// Function overloads via signatures + impl:
function fmt(x: number): string;
function fmt(x: Date): string;
function fmt(x: number | Date): string { ... }
\end{lstlisting}
\end{minipage}
\end{construct}

\begin{construct}{Arrow functions / closures}
\noindent
\begin{minipage}[t]{0.485\linewidth}
\begin{lstlisting}[style=java, title=Java]
Function<Integer,Integer> sq = x -> x * x;
Predicate<String> nonEmpty = s -> !s.isEmpty();
list.forEach(System.out::println);
\end{lstlisting}
\end{minipage}\hfill
\begin{minipage}[t]{0.485\linewidth}
\begin{lstlisting}[style=ts, title=TypeScript]
const sq = (x: number) => x * x;
const nonEmpty = (s: string) => s.length > 0;
xs.forEach((x) => console.log(x));
// Arrow funcs: lexical `this`. Use `function` for own `this`.
\end{lstlisting}
\end{minipage}
\end{construct}

\begin{construct}{Streams / pipelines}
\noindent
\begin{minipage}[t]{0.485\linewidth}
\begin{lstlisting}[style=java, title=Java]
List<Integer> r = xs.stream()
    .filter(x -> x > 0)
    .map(x -> x * 2)
    .sorted()
    .toList();
int total = xs.stream().mapToInt(Integer::intValue).sum();
\end{lstlisting}
\end{minipage}\hfill
\begin{minipage}[t]{0.485\linewidth}
\begin{lstlisting}[style=ts, title=TypeScript]
const r = xs.filter(x => x > 0)
            .map(x => x * 2)
            .sort((a, b) => a - b);
const total = xs.reduce((a, x) => a + x, 0);
\end{lstlisting}
\end{minipage}
\end{construct}

%===============================================================================
\cheatsection{Classes, interfaces, types}
%===============================================================================

\begin{construct}{Class with fields, ctor, methods}
\noindent
\begin{minipage}[t]{0.485\linewidth}
\begin{lstlisting}[style=java, title=Java]
public class Point {
  private final double x, y;
  public Point(double x, double y) {
    this.x = x; this.y = y;
  }
  public double dist(Point o) {
    return Math.hypot(x - o.x, y - o.y);
  }
}
\end{lstlisting}
\end{minipage}\hfill
\begin{minipage}[t]{0.485\linewidth}
\begin{lstlisting}[style=ts, title=TypeScript]
class Point {
  // Param-property shorthand assigns + declares fields.
  constructor(
    public readonly x: number,
    public readonly y: number,
  ) {}
  dist(o: Point): number {
    return Math.hypot(this.x - o.x, this.y - o.y);
  }
}
\end{lstlisting}
\end{minipage}
\end{construct}

\begin{construct}{Inheritance, abstract, interface}
\noindent
\begin{minipage}[t]{0.485\linewidth}
\begin{lstlisting}[style=java, title=Java]
abstract class Shape {
  abstract double area();
}
interface Drawable { void draw(); }
class Circle extends Shape implements Drawable {
  private final double r;
  public Circle(double r) { this.r = r; }
  @Override double area() { return Math.PI*r*r; }
  @Override public void draw() { ... }
}
\end{lstlisting}
\end{minipage}\hfill
\begin{minipage}[t]{0.485\linewidth}
\begin{lstlisting}[style=ts, title=TypeScript]
abstract class Shape {
  abstract area(): number;
}
interface Drawable { draw(): void; }

class Circle extends Shape implements Drawable {
  constructor(private readonly r: number) { super(); }
  area(): number { return Math.PI * this.r ** 2; }
  draw(): void { ... }
}
// Single-class inheritance + multiple interface impl.
\end{lstlisting}
\end{minipage}
\end{construct}

\begin{construct}{Static, getters/setters}
\noindent
\begin{minipage}[t]{0.485\linewidth}
\begin{lstlisting}[style=java, title=Java]
public class Box {
  private int n;
  public int getN() { return n; }
  public void setN(int v) { n = v; }
  public static Box empty() { return new Box(); }
}
\end{lstlisting}
\end{minipage}\hfill
\begin{minipage}[t]{0.485\linewidth}
\begin{lstlisting}[style=ts, title=TypeScript]
class Box {
  private n = 0;
  get value(): number { return this.n; }
  set value(v: number) { this.n = v; }
  static empty(): Box { return new Box(); }
}
const b = new Box(); b.value = 5; console.log(b.value);
\end{lstlisting}
\end{minipage}
\end{construct}

%===============================================================================
\cheatsection{Type system: structural typing}
%===============================================================================

\begin{construct}{Sidebar --- structural vs nominal}
\noindent
\begin{minipage}[t]{\linewidth}
\small Java types are \textbf{nominal}: \code{class A} and \code{class B} with
identical members are unrelated. TypeScript types are \textbf{structural}: shape
matches structure, not name. This is the single biggest mental shift coming from
Java --- it powers a lot of TS's flexibility (and surprises).
\end{minipage}
\end{construct}

\begin{construct}{interface vs type alias}
\noindent
\begin{minipage}[t]{0.485\linewidth}
\begin{lstlisting}[style=java, title=Java]
interface User { String name(); int age(); }
class Member implements User {
  public String name() { return "..."; }
  public int age() { return 0; }
}
// Must declare `implements`.
\end{lstlisting}
\end{minipage}\hfill
\begin{minipage}[t]{0.485\linewidth}
\begin{lstlisting}[style=ts, title=TypeScript]
interface User { name: string; age: number; }
// or:
type User = { name: string; age: number };

const m = { name: "x", age: 1, extra: true };
const u: User = m;          // OK — has all required props

// `interface` is open (mergeable). `type` is closed.
// Use interface for object/class shapes, type for unions.
\end{lstlisting}
\end{minipage}
\end{construct}

\begin{construct}{Unions, literals, narrowing}
\noindent
\begin{minipage}[t]{0.485\linewidth}
\begin{lstlisting}[style=java, title=Java]
// Closest analogs: sealed interfaces, enums.
sealed interface Status permits Active, Banned {}
record Active() implements Status {}
record Banned(String reason) implements Status {}
\end{lstlisting}
\end{minipage}\hfill
\begin{minipage}[t]{0.485\linewidth}
\begin{lstlisting}[style=ts, title=TypeScript]
type Status = "active" | "banned" | "pending";   // literals
type Result<T> = { ok: true; value: T }
              | { ok: false; error: Error };

function handle(r: Result<string>) {
  if (r.ok) {
    r.value;     // narrowed to { ok:true; value:string }
  } else {
    r.error;     // narrowed
  }
}
\end{lstlisting}
\end{minipage}
\end{construct}

\begin{construct}{Type guards \& assertion}
\noindent
\begin{minipage}[t]{0.485\linewidth}
\begin{lstlisting}[style=java, title=Java]
if (s instanceof Member m) { use(m); }
// Pattern match (Java 16+)
\end{lstlisting}
\end{minipage}\hfill
\begin{minipage}[t]{0.485\linewidth}
\begin{lstlisting}[style=ts, title=TypeScript]
if (typeof x === "string") { ... }       // primitive
if (x instanceof Date)     { ... }       // class
if ("foo" in obj)          { ... }       // property test

function isUser(o: unknown): o is User {  // user-defined
  return typeof o === "object" && o !== null
      && "name" in o && typeof (o as any).name === "string";
}

const v = data as User;            // unsafe cast (avoid)
\end{lstlisting}
\end{minipage}
\end{construct}

%===============================================================================
\cheatsection{Generics / type parameters}
%===============================================================================

\begin{construct}{Generic class \& function}
\noindent
\begin{minipage}[t]{0.485\linewidth}
\begin{lstlisting}[style=java, title=Java]
class Box<T> {
  private final T value;
  public Box(T v) { this.value = v; }
  public T get() { return value; }
}
<T extends Comparable<T>> T max(T a, T b) {
    return a.compareTo(b) >= 0 ? a : b;
}
List<? extends Number> ns;
\end{lstlisting}
\end{minipage}\hfill
\begin{minipage}[t]{0.485\linewidth}
\begin{lstlisting}[style=ts, title=TypeScript]
class Box<T> {
  constructor(public readonly value: T) {}
  get(): T { return this.value; }
}
function max<T extends { compareTo(o: T): number }>(
  a: T, b: T,
): T {
  return a.compareTo(b) >= 0 ? a : b;
}
// Bivariance & variance: TS doesn't model wildcards, but
// `readonly T[]` is roughly `List<? extends T>`.
\end{lstlisting}
\end{minipage}
\end{construct}

\begin{construct}{Utility / mapped types}
\noindent
\begin{minipage}[t]{0.485\linewidth}
\begin{lstlisting}[style=java, title=Java]
// No equivalent. You hand-write each variant class.
\end{lstlisting}
\end{minipage}\hfill
\begin{minipage}[t]{0.485\linewidth}
\begin{lstlisting}[style=ts, title=TypeScript]
type User = { name: string; age: number; email: string };
type PartialUser = Partial<User>;       // all optional
type ReadOnlyUser = Readonly<User>;
type Email       = Pick<User, "email">;
type NoEmail     = Omit<User, "email">;
type RequiredU   = Required<PartialUser>;
type Keys        = keyof User;          // "name"|"age"|...
// Mapped types compile to runtime-erased types.
\end{lstlisting}
\end{minipage}
\end{construct}

%===============================================================================
\cheatsection{Error handling}
%===============================================================================

\begin{construct}{Try / catch / finally}
\noindent
\begin{minipage}[t]{0.485\linewidth}
\begin{lstlisting}[style=java, title=Java]
try (var f = Files.newBufferedReader(p)) {
    return f.readLine();
} catch (IOException e) {
    log.error("read failed", e);
    throw new MyException("oops", e);
} finally { cleanup(); }
\end{lstlisting}
\end{minipage}\hfill
\begin{minipage}[t]{0.485\linewidth}
\begin{lstlisting}[style=ts, title=TypeScript]
try {
    return await readFile(p, "utf8");
} catch (e) {
    // Caught value type is `unknown` (TS 4.4+).
    if (e instanceof Error) {
        log.error(`read failed: ${e.message}`);
    }
    throw new MyError("oops", { cause: e });
} finally { cleanup(); }
\end{lstlisting}
\end{minipage}
\end{construct}

\begin{construct}{Custom error}
\noindent
\begin{minipage}[t]{0.485\linewidth}
\begin{lstlisting}[style=java, title=Java]
public class NotFoundException extends RuntimeException {
  public NotFoundException(String m) { super(m); }
}
throw new NotFoundException("user " + id);
\end{lstlisting}
\end{minipage}\hfill
\begin{minipage}[t]{0.485\linewidth}
\begin{lstlisting}[style=ts, title=TypeScript]
class NotFoundError extends Error {
  constructor(message: string) {
    super(message);
    this.name = "NotFoundError";
  }
}
throw new NotFoundError(`user ${id}`);
\end{lstlisting}
\end{minipage}
\end{construct}

%===============================================================================
\cheatsection{Modules / packages / imports}
%===============================================================================

\begin{construct}{Imports / exports}
\noindent
\begin{minipage}[t]{0.485\linewidth}
\begin{lstlisting}[style=java, title=Java]
package com.acme.svc;
import java.util.List;
import static com.acme.Utils.foo;
\end{lstlisting}
\end{minipage}\hfill
\begin{minipage}[t]{0.485\linewidth}
\begin{lstlisting}[style=ts, title=TypeScript]
import { foo, bar } from "./utils";              // named
import defaultExport from "./mod";               // default
import * as utils from "./utils";                // namespace
import type { User } from "./types";             // type-only

export const x = 1;
export function f() { ... }
export default class C { ... }
export type { User };
// "./utils.js" extension required in NodeNext / ESM mode.
\end{lstlisting}
\end{minipage}
\end{construct}

%===============================================================================
\cheatsection{Concurrency: Promises \& async/await}
%===============================================================================

\begin{construct}{Async functions}
\noindent
\begin{minipage}[t]{0.485\linewidth}
\begin{lstlisting}[style=java, title=Java]
CompletableFuture<String> fut =
  CompletableFuture.supplyAsync(() -> fetch())
    .thenApply(String::trim)
    .exceptionally(e -> "fallback");
String r = fut.get();
\end{lstlisting}
\end{minipage}\hfill
\begin{minipage}[t]{0.485\linewidth}
\begin{lstlisting}[style=ts, title=TypeScript]
async function run(): Promise<string> {
    try {
        const raw = await fetchData();
        return raw.trim();
    } catch {
        return "fallback";
    }
}
const r = await run();        // top-level await in modules
\end{lstlisting}
\end{minipage}
\end{construct}

\begin{construct}{Combinators \& parallel}
\noindent
\begin{minipage}[t]{0.485\linewidth}
\begin{lstlisting}[style=java, title=Java]
CompletableFuture<Void> all = CompletableFuture.allOf(f1, f2);
CompletableFuture<Object> any = CompletableFuture.anyOf(f1, f2);
\end{lstlisting}
\end{minipage}\hfill
\begin{minipage}[t]{0.485\linewidth}
\begin{lstlisting}[style=ts, title=TypeScript]
const [a, b] = await Promise.all([p1, p2]);     // all OK
const r = await Promise.race([p1, p2]);
const settled = await Promise.allSettled([p1,p2]);
const first  = await Promise.any([p1, p2]);     // any FULFILL
\end{lstlisting}
\end{minipage}
\end{construct}

\begin{construct}{Single-threaded event loop}
\noindent
\begin{minipage}[t]{\linewidth}
\small
TypeScript runs on JS engines, which are \textbf{single-threaded} by default.
\code{async}/\code{await} schedules continuations on the event loop --- no
preemption, no shared-memory races on JS objects. For CPU-heavy work, spawn
\code{Worker} threads (browser) or \code{worker_threads} (Node), which have
isolated memory and communicate via \code{postMessage}. This is a fundamental
shift from Java's pre\"emptive threading + shared mutable state model.
\end{minipage}
\end{construct}

%===============================================================================
\cheatsection{Idiom appendix}
%===============================================================================

\begin{construct}{Equality, null, undefined}
\noindent
\begin{minipage}[t]{0.485\linewidth}
\begin{lstlisting}[style=java, title=Java]
if (s == null) { ... }              // identity
if (a.equals(b)) { ... }            // value
String x = (s != null) ? s : "default";
\end{lstlisting}
\end{minipage}\hfill
\begin{minipage}[t]{0.485\linewidth}
\begin{lstlisting}[style=ts, title=TypeScript]
if (s === null) { ... }             // strict equality
if (s == null)  { ... }             // null OR undefined
// Always prefer === / !== — `==` does coercion.
const x = s ?? "default";           // null/undefined only
const y = s || "default";           // ALSO falsy: "", 0, false
\end{lstlisting}
\end{minipage}
\end{construct}

\begin{construct}{Common gotchas (TS coming from Java)}
\noindent
\begin{minipage}[t]{\linewidth}
\small
\begin{itemize}[leftmargin=*, itemsep=1pt, topsep=1pt]
  \item \textbf{Types erase at runtime.} \code{instanceof T} works only for classes; you can't \code{instanceof MyInterface}. Use type guards or a \code{kind} discriminator.
  \item \textbf{this binding.} Arrow functions capture \code{this} lexically; \code{function} expressions don't. Methods passed as callbacks lose \code{this} unless bound.
  \item \textbf{undefined vs null.} Both exist. Default uninitialized = \code{undefined}. Prefer one consistently; many codebases ban \code{null}.
  \item \textbf{== is coercion.} \code{0 == ""} is true. Always \code{===}/\code{!==}.
  \item \textbf{Array holes.} \code{[1,,3]} has a sparse slot; \code{map} skips it. Avoid.
  \item \textbf{any vs unknown.} \code{any} disables type-checking. \code{unknown} forces narrowing before use --- prefer it.
  \item \textbf{strict mode.} Always set \code{"strict": true} in \code{tsconfig.json}; otherwise nullability checks are off.
  \item \textbf{Numbers are floats.} No int. \code{0.1 + 0.2 !== 0.3}. Use \code{bigint} for exact integers, libraries for decimal money math.
  \item \textbf{Tooling.} \code{tsc} (compiler), \code{tsx}/\code{ts-node} (run), \code{eslint}, \code{prettier}, \code{vitest}/\code{jest}, \code{tsd}/\code{expect-type} for type-level tests.
\end{itemize}
\end{minipage}
\end{construct}

\end{document}
